% =====================================================================
\chapter{Cuestiones}
\label{cap:test}
% =====================================================================

Veinte cuestiones de opción múltiple, de 2 puntos cada una (40 puntos en total). Hay una sola respuesta correcta en cada cuestión y las respuestas erróneas no restan. La clave y su justificación están en el \cref{ap:soluciones}.

\begin{cuestiones}
\item \pt{2} Número de FLOP necesarios para entrenar un modelo de $3\cdot 10^{9}$ parámetros con $6\cdot 10^{10}$ tokens:
\begin{opciones}
  \item $1{,}8\cdot 10^{20}$
  \item $3{,}6\cdot 10^{20}$
  \item $1{,}08\cdot 10^{21}$
  \item $1{,}08\cdot 10^{22}$
\end{opciones}

\item \pt{2} Según la regla aproximada de Chinchilla de unos 20 tokens por parámetro, ¿cuántos tokens corresponden a un modelo de $8\cdot 10^{9}$ parámetros?
\begin{opciones}
  \item $1{,}6\cdot 10^{11}$
  \item $8\cdot 10^{9}$
  \item $4\cdot 10^{10}$
  \item $1{,}5\cdot 10^{13}$
\end{opciones}

\item \pt{2} Frente a Kaplan \emph{et al.} (2020), Hoffmann \emph{et al.} (Chinchilla, 2022) concluyeron que, para un cómputo fijo, conviene entrenar:
\begin{opciones}
  \item modelos más grandes con menos datos;
  \item con el mismo reparto entre parámetros y datos;
  \item más épocas sobre los mismos datos;
  \item modelos más pequeños con más datos.
\end{opciones}

\item \pt{2} Un modelo de $1{,}3\cdot 10^{10}$ parámetros se entrena con Adam, precisión mixta y ZeRO-3 en 16 GPU. Memoria de estado por GPU, sin contar las activaciones:
\begin{opciones}
  \item 1{,}6~GB
  \item 13~GB
  \item 26~GB
  \item 208~GB
\end{opciones}

\item \pt{2} Intensidad aritmética crítica de la H100 (989~TFLOP/s; 3{,}35~TB/s), en FLOP/byte:
\begin{opciones}
  \item $\approx 30$
  \item $\approx 156$
  \item $\approx 295$
  \item $\approx 3300$
\end{opciones}

\item \pt{2} Cota superior de la velocidad de decodificación, con lote~1, de un modelo de $3\cdot 10^{10}$ parámetros en bf16 en una H100, ignorando la caché KV:
\begin{opciones}
  \item $\approx 56$ tokens/s
  \item $\approx 28$ tokens/s
  \item $\approx 112$ tokens/s
  \item $\approx 16\,500$ tokens/s
\end{opciones}

\item \pt{2} Caché KV por token de un modelo con $L=40$, $\nkv=8$ y $\dhd=128$, en bf16:
\begin{opciones}
  \item 40~KiB
  \item 80~KiB
  \item 1{,}25~MiB
  \item 160~KiB
\end{opciones}

\item \pt{2} Fracción del tiempo ociosa (burbuja) de un pipeline de GPipe con $p_{\mathrm{e}}=8$ etapas y $m=24$ microlotes:
\begin{opciones}
  \item 7{,}0\,\%
  \item 22{,}6\,\%
  \item 25{,}0\,\%
  \item 33{,}3\,\%
\end{opciones}

\item \pt{2} Grok-1, cuyos pesos se publicaron en 2024, tiene unos $3{,}14\cdot 10^{11}$ parámetros, y xAI indicó que en cada token se activa alrededor del 25\,\% de los pesos. Parámetros activos por token, aproximadamente:
\begin{opciones}
  \item $7{,}85\cdot 10^{10}$
  \item $3{,}9\cdot 10^{10}$
  \item $1{,}57\cdot 10^{11}$
  \item $3{,}14\cdot 10^{11}$
\end{opciones}

\item \pt{2} En una capa de mezcla de expertos con enrutado top-2, al pasar de 8 a 64 expertos del mismo tamaño crece:
\begin{opciones}
  \item el número de FLOP por token;
  \item la caché KV;
  \item la memoria de pesos;
  \item la longitud de contexto.
\end{opciones}

\item \pt{2} La principal ventaja de la atención de consultas agrupadas (GQA) frente a la atención multicabeza clásica es:
\begin{opciones}
  \item menos FLOP en el MLP;
  \item un contexto ilimitado;
  \item entrenar sin máscara causal;
  \item una caché KV más pequeña.
\end{opciones}

\item \pt{2} Atención con ventana deslizante $W=2048$ en 24 capas apiladas. Campo receptivo teórico aproximado, $L\cdot W$:
\begin{opciones}
  \item 2048 tokens
  \item $49\,152$ tokens
  \item $24\,576$ tokens
  \item ilimitado
\end{opciones}

\item \pt{2} Cuatro expertos con fracciones de tráfico $f=(0{,}4;\ 0{,}3;\ 0{,}2;\ 0{,}1)$ y probabilidades medias del enrutador $P=(0{,}25;\ 0{,}25;\ 0{,}25;\ 0{,}25)$. Valor de $\Lbal=N\sum_i f_iP_i$:
\begin{opciones}
  \item $1$
  \item $0{,}25$
  \item $1{,}2$
  \item $4$
\end{opciones}

\item \pt{2} ¿Qué afirmación sobre FlashAttention es correcta? ($n$ es la longitud de la secuencia.)
\begin{opciones}
  \item Aproxima la softmax para ir más rápido.
  \item Reduce los FLOP de la atención a $O(n)$.
  \item Es exacta y su memoria adicional en HBM es $O(n)$.
  \item Solo puede usarse en inferencia.
\end{opciones}

\item \pt{2} El \emph{prefill} de un prompt largo está limitado sobre todo por:
\begin{opciones}
  \item el ancho de banda de memoria;
  \item la red entre nodos;
  \item el tokenizador;
  \item la capacidad de cómputo.
\end{opciones}

\item \pt{2} La decodificación especulativa con la regla de aceptación $\min(1,p/q)$ y remuestreo del residuo:
\begin{opciones}
  \item aproxima la distribución $p$ del modelo objetivo con un error pequeño;
  \item reproduce exactamente la distribución $p$ del modelo objetivo;
  \item solo funciona con decodificación voraz;
  \item reduce la memoria de pesos del modelo objetivo.
\end{opciones}

\item \pt{2} Bits por peso en int4 con una escala fp16 por cada grupo de 64 pesos:
\begin{opciones}
  \item $4{,}25$
  \item $4$
  \item $4{,}125$
  \item $4{,}5$
\end{opciones}

\item \pt{2} Una respuesta de 20 tokens solo cuenta como correcta si acierta todos, y cada token se acierta, de forma independiente, con probabilidad $0{,}95$. La exactitud exacta vale, aproximadamente:
\begin{opciones}
  \item $0{,}95$
  \item $0{,}60$
  \item $0{,}36$
  \item $0{,}05$
\end{opciones}

\item \pt{2} PagedAttention (vLLM) aumenta el rendimiento del servicio porque:
\begin{opciones}
  \item reduce los FLOP de la atención;
  \item cuantiza la caché KV;
  \item elimina la fase de \emph{prefill};
  \item evita reservar para cada petición una caché KV contigua del tamaño máximo.
\end{opciones}

\item \pt{2} La utilización de FLOP del modelo (MFU) de un entrenamiento es:
\begin{opciones}
  \item la memoria usada dividida por la memoria total;
  \item los FLOP útiles del modelo por segundo divididos por el rendimiento máximo teórico;
  \item los tokens por segundo y por GPU;
  \item los vatios consumidos por FLOP.
\end{opciones}
\end{cuestiones}
